\subsection{Setting}

In this paper, we introduce and study a refinement of the relative (or logarithmic) Gromov--Witten invariants \cite{abramovich2014stable,gross2013logarithmic,kim2010logarithmic,li2001stable} of the pair $(Y=\PP_X(\O\oplus L),D)$, where $D=D^+ + D^-=\PP_X(L)\oplus\PP_X(\O)$ and $X$ is a smooth projective variety over $\mathbb C$. We fix the following discrete data:
\begin{itemize}\itemsep=0pt
 \item $g$, $n$, $m$ are positive integers,
 \item $\bfw=(w_1,\dots,w_n)$ is a $n$-tuple of non-zero integers with $\sum w_i=0$; we moreover set \smash{$b(\bfw)=\sum_{w_i>0}w_i=-\sum_{w_i<0}w_i$};
 \item $\beta\in H_2(X,\ZZ)$ is an effective homology class.
\end{itemize}
Let $\M:=\M_{g,m}\bigl(Y|D^{\pm},\beta,\bfw\bigr)$ be the moduli space of logarithmic stable maps parametrizing
\[
f\colon\ \ (C,p_1,\dots,p_m,q_1,\dots,q_n)\to Y,
 \]
where $C$ is a stable curve of genus $g$ with image in class $\beta+b(\bfw)F\in H_2(Y,\ZZ)$, the marked point $q_i$ is mapped to $D^{\mathrm{sgn}(w_i)}$ with tangency order $|w_i|$, the marked points $p_i$ are mapped to the interior. The vector $\bfw$ is thus called \textit{tangency profile}.
The logarithmic (or equivalently relative~\cite{abramovich2014comparison}) Gromov--Witten invariants are defined integrating constraints pulled-back along evaluation maps~${
\ev \colon \M_{g,m}\bigl(Y|D^{\pm},\beta,\bfw\bigr) \longrightarrow X^n\times Y^m}$,
 against the virtual fundamental class~\smash{$\vir{\M}$}~\cite{gross2013logarithmic}.

 We propose a refinement for these GW invariants by remembering a discrete quantity depending on the position of the points mapped to the boundary divisor. This quantity is called a \textit{correlator}.

\subsection{The Albanese evaluation}
The idea of the refinement comes from reinterpreting the map to $\bigl(Y, D^\pm\bigr)=\bigl(\PP_X(\O\oplus L), D^\pm\bigr)$ as the data of a map $f\colon C\to X$ together with the choice of an isomorphism $f^*L\cong\O_C(\alpha)$ where $\alpha\in H^0\bigl(C,\bar{M}_C^{\rm gp}\bigr)$ is defined using the logarithmic structure on $C$. When $C$ is smooth, this condition takes the familiar form $f^*L\simeq\O_C(\sum w_i q_i)$.

\subsubsection[Case of X times P\^{}1]{Case of $\boldsymbol{ X\times\PP^1}$}


To simplify, let us momentarily assume that $L$ is the trivial bundle and that the source curve~$C$ is smooth.
 Then, reinterpreting $f\in\M$ as explained above, part of the data of a log stable map is the isomorphism
$
\O_C\bigl(\sum w_iq_i\bigr)\cong\O\in\Pic^0(C)$.
For a smooth curve, the Abel--Jacobi theorem ensures that $\Pic^0(C)\simeq\Alb(C)$, where $\Alb(C)$ is the Albanese variety of $C$.
 We denote by~${a_X\colon X\to\Alb(X)}$ the map from $X$ to its Albanese variety, unique up to translation. By the functoriality of the Albanese variety, $f$ induces a~morphism of Abelian varieties
$
 f_*\colon \Alb(C)\to\Alb(X)
$ compatible with the Albanese maps.

We denote by $a_\bfw$ the morphism
\[
 a_\bfw\colon\ (x_1,\dots,x_n)\in X^n\longmapsto \sum_{i=1}^n w_i a_X(x_i)\in\Alb(X).
 \]
As $\sum w_i=0$, the map $a_\bfw$ does not change if we compose $a_X$ with a translation. From $\O_C(\sum w_iq_i)\simeq\O\in\Pic^0(C)$, we get that
\begin{equation}\label{eq-f*f*}
 \sum_{i=1}^n w_i a_X(f(q_i)) =0\in\Alb(X).
\end{equation}
Equation \eqref{eq-f*f*} is a constraint on the relative position of the images $f(q_i)$. In other words, the~$X^n$ part of the evaluation map is not surjective, and truly has values in the subset $a_\bfw^{-1}(0)$.


\subsubsection{Case of a non-trivial bundle} If $L$ is not assumed to be trivial anymore, we prove in Lemma~\ref{lem:mapphibeta} that there is a natural map $\varphi_\beta\colon \Pic^0(X)\to\Alb(X)$ defined through Hodge theory, which only depend on the curve class. Then
Proposition \ref{prop-image-of-L-is-constant} ensures that equation \eqref{eq-f*f*} becomes the following:
\begin{equation}\label{eq-f*f*L}
\sum_{i=1}^n w_ia_X(f(q_i)) = \varphi_\beta(L)\in\Alb(X).
\end{equation}
In other words, the image depends only on the degree $\beta$ and the line bundle $L$ and not on $f$.

\subsection{Correlated classes and correlated GW invariants} We now assume that the tangency orders $w_i$ have a non-trivial common divisor $\delta$. We then consider the map
\[
a_{\bfw/\delta}\colon\ (x_1,\dots,x_n)\in X^n\longmapsto\sum_1^n\frac{w_i}{\delta}a_X(x_i)\in \Alb(X).
 \]
The above morphism induces a morphism from the space of log stable maps $\M$
\[\kappa^\delta(f)=\sum \frac{w_i}{\delta}a_X(f(q_i))\in\Alb(X).\]
It follows from equation \eqref{eq-f*f*L} that
$\delta\cdot\kappa^\delta(f)=\varphi_\beta(L)$.
Therefore, the map $\kappa^\delta$ takes values in the set of $\delta$-roots of $\varphi_\beta(L)$, denoted by $T^\delta(L,\beta)$. The latter is a torsor under the group of $\delta$-torsion elements in $\Alb(X)$, denoted by $\Tor_\delta(\Alb(X))$.
The value of $\kappa^\delta$ is called a \textit{correlator} and is denoted by $\theta$.
We want to refine the Gromov--Witten of $\bigl(Y,D^{\pm}\bigr)$ keeping track of the value of the correlator.

The discussion above assumes the source curve $C$ to be smooth, but the morphism $\kappa^\delta$
is also defined on the boundary of the moduli space $\M_{g,m}\bigl(Y|D^{\pm},\beta,\bfw\bigr)$, still with values in $T^\delta(L,\beta)$ (see Section~\ref{sec:correlationrefinement}).
This shows that
 $\M$ splits into distinct (possibly still disconnected) components~$\M^\theta$ indexed by the values of the correlator $\theta\in T^\delta(L,\beta)$; in particular the virtual class splits as a sum of the so-called \textit{correlated classes} \smash{$\vir{\M^\theta}$}. We define the \textit{full correlated class}~$\fvir{\M}{\delta}$
\[ \fvir{\M}{\delta}=\sum_{\theta\in T^\delta(L,\beta)} \vir{\M^\theta}\cdot(\theta)\in \QQ[\Alb(X)]\otimes H_\bullet(\M,\QQ), \]
which is an element of the group algebra with cycle coefficients with support on $T^\delta(L,\beta)\subset\Alb(X)$.

\textit{Correlated GW invariants} are obtained by capping these classes with some pullback by the evaluation map: if $\gamma\in H^\bullet(X^n\times Y^m,\QQ)$, we set
\[ \langle\langle\gamma\rangle\rangle^\delta_{g,\beta,\bfw} = \int_{\fvir{\M}{\delta}}\ev^*\gamma \in\QQ[\Alb(X)], \]
which is an element of the group algebra $\QQ[\Alb(X)]$ with support on the torsor $T^\delta(L,\beta)$. Distinct correlators may yield different Gromov--Witten invariants, as illustrated by the computations in the elliptic case.

It follows from the discussion above that the evaluation map truly has values in the set \smash{$a_\bfw^{-1}(\varphi_\beta(L))=\bigsqcup a_{\bfw/\delta}^{-1}(\theta)$}. A broader generalization of GW invariants would be obtained by allowing the pullback of cohomology classes from $a_\bfw^{-1}(\varphi_\beta(L))$ rather than $X^n$. Correlated invariants are a particular case of these invariants where we pull-back the classes $1_\theta\in H^0\bigl(a_\bfw^{-1}(\varphi_\beta(L)),\QQ\bigr)$, corresponding to the components $a_{\bfw/\delta}^{-1}(\theta)$ of $a_{\bfw}^{-1}(\varphi_\beta(L))$.


